\subsection{The ground line with varying physical dimensions}

The adjacent open parent has a unique normalized ground-state line even when
its physical dimensions vary from site to site. The length condition $N\ge2$
is retained explicitly; see \cite{gap:pgvwc07_open_parent_domain}.

\begin{definition}[Common-alphabet extension of the site families]
    \label{def:varying_physical_actual_site_extension}
    \lean{Matrix.commonPhysicalSiteFamily, Matrix.commonPhysicalSiteFamily_apply}
    \leanok
    \uses{def:varying_physical_common_alphabet}
    For $0\le j<N$, let $B^{[j]}_i$ be $D_j\times D_{j+1}$ matrices with
    $0\le i<d_j$. Put $m=\sum_{j<N}d_j$ and choose the canonical inclusions
    $e_j:\{0,\ldots,d_j-1\}\hookrightarrow\{0,\ldots,m-1\}$.
    The common-alphabet families are
    \begin{align}
        C^{[j]}_{e_j(i)} & =B^{[j]}_i,                                   &
        C^{[j]}_t        & =0\quad\text{if }t\notin\operatorname{im}e_j.
    \end{align}
    Families outside the finite chain are zero.
\end{definition}

\begin{definition}[The state on the original physical alphabets]
    \label{def:varying_physical_scalar_state}
    \lean{Matrix.varyingPhysicalScalarState}
    \leanok
    \uses{def:varying_physical_actual_site_extension, def:varying_obc_scalar_state}
    Suppose $D_0=D_N=1$. Read the common-alphabet product at the
    embedded configurations. This gives the vector $\Omega$ on the original
    physical configuration space with coefficients
    \begin{align}
        \Omega(i_0,\ldots,i_{N-1})
         & =\bigl(C^{[0]}_{e_0(i_0)}\cdots
        C^{[N-1]}_{e_{N-1}(i_{N-1})}\bigr)_{0,0}                       \\
         & =\bigl(B^{[0]}_{i_0}\cdots B^{[N-1]}_{i_{N-1}}\bigr)_{0,0}.
    \end{align}
\end{definition}

\begin{lemma}[Normalization on the original physical space]
    \label{lem:varying_physical_scalar_state_normalized}
    \lean{Matrix.norm_varyingPhysicalScalarState_eq_one}
    \leanok
    \uses{def:varying_physical_scalar_state,
        lem:varying_physical_kernel_actual_pullback,
        lem:uniform_physical_parent_kernel_comparison,
        lem:varying_obc_scalar_state_normalized}
    Let $N\ge2$, $D_0=D_N=1$, and suppose that
    $\sum_iB^{[j]}_iB^{[j]\dagger}_i=I_{D_j}$ at each site. Then $\|\Omega\|=1$.
\end{lemma}
\begin{proof}\leanok
    Extension by zero preserves the canonical row equations, so the common-alphabet state $\widetilde\Omega$ has norm one. In particular the
    common alphabet is nonempty. Its parent equation forces
    $\widetilde\Omega=E\Omega$, where $E$ is the isometric physical extension.
    Hence $\|\Omega\|=\|E\Omega\|=1$.
\end{proof}

\begin{lemma}[The ground line for arbitrary physical dimensions]
    \label{lem:varying_physical_actual_parent_unique_line}
    \lean{Matrix.ker_varyingPhysicalParentHamiltonian_eq_span,
        Matrix.finrank_ker_varyingPhysicalParentHamiltonian_eq_one}
    \leanok
    \uses{lem:varying_physical_scalar_state_normalized,
        lem:uniform_physical_parent_kernel_comparison,
        lem:varying_obc_actual_parent_unique_line,
        lem:varying_physical_parent_embedding_transport}
    Consider $N\ge2$ sites with arbitrary physical dimensions $d_j$, singleton
    endpoint bonds, and canonical row equations. Suppose each site family
    spans its full rectangular matrix space. Let $H$ be the sum of the
    adjacent complementary pair projections on the original physical space.
    Then
    \begin{align}
        \ker H & =\mathbb C\Omega, & \dim\ker H & =1.
    \end{align}
    Since $H\ge0$, this is its unique ground-state line.
\end{lemma}
\begin{proof}\leanok
    The extended site families retain the canonical row equations and their
    full rectangular spans. Parent transport and the common-alphabet
    comparison carry every vector $v\in\ker H$ to a vector $Ev=cE\Omega$ in
    the common-alphabet ground line. Injectivity gives $v=c\Omega$.
    Conversely $E\Omega$ is annihilated by that parent, and the derived
    transport identity gives $H\Omega=0$. The unit norm proves that the
    resulting kernel line is nonzero and one-dimensional.
\end{proof}
